\section{Details of the computation}\label{app:computation}
The invariants $\brCS_\lambda$ and $\hf_\rho$ are evaluated modulo a $31$-bit prime $p$ on a grid and interpolated in each of the variables $(A,Q,T)=(a,q,t)^{1/2}$.
When coefficients exceed $p/2$, two or three primes are used and the result is reconstructed by the Chinese remainder theorem (Garner's algorithm).
The polynomials obtained are checked at points not used in the interpolation.
Evaluation points are stored in blocks so that a computation can be resumed.

For the genus expansion we proceed as in Lemma~\ref{lem:G}: with $v=\bq^2\bt$, $w=v+v^{-1}$ and $w_0=\bt+\bt^{-1}$ we have $G=w-w_0$.
Each coefficient of a power of $\ba$ is rewritten as a polynomial in $w$ and expanded in powers of $G$ by Horner's scheme.

For $|\rho|\ge4$, $\hf_\rho$ is computed directly at each evaluation point (\S\ref{sec:higher}).
The variable $A$ is also evaluated numerically and interpolated as a polynomial in $a=A^2$, with a check at one point outside the range.
For torus knots the parities of the exponents of $Q$ and $T$ are uniform, so we interpolate in $q=Q^2$, $t=T^2$, which reduces the number of evaluation points by a factor of four.
In the elliptic Hall algebra computation, the action of the operators corresponding to the power sums and its evaluation are shared across the colours $\lambda$.
The running times are those of Table~\ref{tab:higher}, with eight parallel processes.

The computations were done on an Apple M3 Pro with $36$\,GB of memory and a Mac mini M1 with $8$\,GB.
The heaviest case is $K10$: $\rCS_{\yng(3)}$ took about four hours ($14440$\,s) with two primes and $\rCS_{\yng(2,1)}$ about two hours ($7387$\,s).
The three-box $\hf_\rho$ required three primes, $4050$\,s and at most $17.6$\,GB of memory.
For $K7$ and $K8$ all colours with $|\rho|\le3$ took $12275$\,s and $22853$\,s, with less than $2$\,GB of memory.

Each stage of the implementation was tested on small inputs:
the reproduction of the values of KN, of the explicit examples of Cherednik and Danilenko and of HOMFLY-PT polynomials computed from braids,
the interpolation and the Chinese remainder theorem, and the agreement of two independent implementations of the genus expansion.

\section{Derivation of the normalisation}\label{app:normalization}
We give the details of the proof of Proposition~\ref{thm:M}, first for the cabling step and then for torus knots.
The framing operator $\tau:P_\mu\mapsto T_\mu P_\mu$ satisfies $\tau^jP_\lambda[p_k\mapsto P_{k(m,n)}]\cdot1=P_\lambda[p_k\mapsto P_{k(m,n+jm)}]\cdot1$ (Computational Result~\ref{cr:tau}).
The cable construction is linear in the companion. Applying (N2$'$) to the companion shifts the slope of the pattern from $(m,-s')$ to $(m,-n)$, i.e.\ to the Seifert framing.
Together with Hypothesis~\ref{hyp:sat} this gives $M_{C_{m,n}K}=M_K^{\,m}M_{T_{m,n}}$ and $F=mn$.

For torus knots the four facts used in the proof are the following.
(I) The coefficient of $a^0$ in the unnormalised reduced value of Gorsky and Negu\c{t} is $(-t/q)^{m-1}C_{m,n}(1/q,t)$, where $C_{m,n}$ is the rational $q,t$-Catalan number.
This was conjectured by Gorsky and Negu\c{t} and follows from the theorem of Mellit \cite{Mellit16} through the operators $Q_{m,n}$ of Bergeron, Garsia, Leven and Xin.
(II) The lowest $a$-degree $\delta$ follows from the theorem of Franks and Williams on positive braids.
(III) The sign and the power of $q$ at $t=q$ follow from the $q,t$-symmetry of $C_{m,n}$.
(IV) The exponent of $(q/t)$ follows from the Dyck-path form of Mellit \cite{Mellit17} and Computational Result~\ref{cr:mellit}.
Together these give $M_{T_{m,n}}=(-q/t)^{m-1}\bigl(a(q/t)^{1/2}\bigr)^{\delta}$.

The formula of Computational Result~\ref{cr:mellit} used in (IV) differs in form from the corollary as stated by Mellit.
Mellit \cite[Cor.~3.4]{Mellit17} writes $\mathcal{P}_{m,n}=(-1)^n(a(qt)^{-1/2})^\delta\frac{1}{1-t}\frac{P_{m,n}(1)[a-1]}{a-1}$ with $P_{m,n}$ ``the corresponding generator of the elliptic Hall algebra''.
By Thm.~3.3 of the same paper this holds if $P_{m,n}(1)$ is read as $-d_-B_{c_{m,n}}d_+(1)$ in terms of the operators $d_\pm$ of Carlsson and Mellit.
If instead it is read as the element in the standard normalisation ($Q_{m,n}$ of \cite{Bergeron14} or $\widetilde P$ of \cite{Gorsky13b}), the formula with the substitution $[a-1]$ does not hold;
among the variants we tried, only the form of Computational Result~\ref{cr:mellit} holds.
Since (IV) only fixes the exponent of $(q/t)$ in $M_\square$, this point does not affect any positivity result (Lemma~\ref{lem:adams}).

\section{Non-algebraic cables}\label{app:nonalg}
The non-algebraic cables that we computed are the $(2,1)$-, $(2,3)$-, $(2,5)$-, $(2,7)$- and $(2,-1)$-cables of $T_{2,3}$, named $K2$, $K4$, $K5a$, $K5b$, $K5c$ in this order,
and the $(2,1)$-cable $K6b$ of $T_{2,5}$.
These knots have $n_2<m_1n_1m_2$. For the five cables other than $K5b$ the DAHA value for $\yng(1)$ does not satisfy the necessary conditions coming from spectral sequences and is therefore not the Poincar\'e polynomial of HOMFLY-PT homology,
so the premise of the KN positivity is not satisfied. Table~\ref{tab:nonalg} shows the outcome of the tests.
\begin{table}[h]
\centering
\caption{Non-algebraic cables at framing zero. Each entry reads ``negative terms/terms of the spin form; negative terms/terms of the genus form'', with the better overall sign in each case.}\label{tab:nonalg}
\sbox\tabbox{\begin{tabular}{llcccccc}
\hline
Knot & Class & $\yng(1)$ & $\yng(2)$ & $\yng(1,1)$ & $\yng(3)$ & $\yng(2,1)$ & $\yng(1,1,1)$ \\
\hline
$K2$ & non-alg. & 10/47; 6/48 & 152/1136; 41/1045 & 174/1273; 51/1190 & 967/8525; 253/7852 & 1078/9527; 292/8774 & 1156/9990; 318/9254 \\
$K4$ & non-alg. & 18/42; 20/43 & 286/1109; 195/952 & 330/1255; 224/1090 & 2561/8614; 1749/7420 & 2850/9646; 1989/8325 & 3025/10121; 2128/8795 \\
$K5a$ & non-alg. & 8/38; 5/47 & 409/1226; 332/1066 & 460/1377; 376/1210 & 4441/9758; 3631/8278 & 4982/10844; 4066/9235 & 5164/11348; 4263/9733 \\
$K5b$ & non-alg. & \checkmark\,(48) & 136/1395; 66/1305 & 149/1537; 78/1452 & 2097/10706; 1060/9612 & 2358/11707; 1181/10553 & 2351/12202; 1228/11040 \\
$K5c$ & non-alg. & 5/57; 0/69 & 78/1307; 0/1299 & 77/1443; 0/1433 & 400/9501; 6/9123 & 423/10494; 7/10053 & 376/10864; 0/10444 \\
$K6b$ & non-alg. & 30/133; 10/178 & 576/4400; 81/4367 & 618/4677; 91/4656 & 3626/33214; 512/32175 & 3766/35136; 541/34032 & 3878/36067; 560/34980 \\
\hline
\end{tabular}
}\resizebox{\textwidth}{!}{\usebox\tabbox}
\end{table}
The failures are of three types.
(I) For the $(2,3)$- and $(2,5)$-cables of $T_{2,3}$ ($K4$, $K5a$), positivity with a refinement of uniform parity is impossible already at the level of the uncoloured HOMFLY-PT polynomial.
(II) For the $(2,1)$-cables ($K2$, $K6b$) and the $(2,-1)$-cable ($K5c$) the HOMFLY-PT polynomial is compatible with positivity, but the DAHA value $\rCS_{\yng(1)}$ itself has negative coefficients.
(III) For the $(2,7)$-cable of $T_{2,3}$ ($K5b$) the failure in the coloured case depends on the framing:
among even framings $|x|\le24$, both forms are positive for two boxes when $x\le-12$ and for three boxes when $x\le-14$.
The cable $K5c$ is an example in which the genus form is (almost) non-negative while the spin form is not, showing that the converse of Lemma~\ref{lem:G} fails.

\section{Data and code}\label{app:data}
The values of $\brCS_\lambda$, $\hf_\rho$, the spin form and the genus form for all cases reported here,
the certificates used in the computer-assisted proofs, and the code and tests used for the computation will be made available with the final version of this paper.
